% !TeX program = pdflatex
% Compilar en Overleaf con pdfLaTeX. Las figuras se incluyen en la carpeta figuras/.
% Proyecto de Sistemas Inteligentes - FES Acatlán, UNAM
\documentclass[aspectratio=169,11pt]{beamer}
\usepackage[utf8]{inputenc}
\usepackage[T1]{fontenc}
\usepackage[spanish,es-noquoting]{babel}
\usepackage{lmodern}
\usepackage{amsmath,amssymb}
\usepackage{graphicx}
\usepackage{booktabs}
\usepackage{array}
\usepackage{tikz}
\usetikzlibrary{arrows.meta,positioning,shapes.geometric,calc}
\usepackage{pgfplots}
\pgfplotsset{compat=1.18}
\usepackage{xcolor}
\usepackage{listings}
\usepackage{hyperref}

\definecolor{UNAMnavy}{HTML}{14284F}
\definecolor{UNAMdeep}{HTML}{0C1730}
\definecolor{UNAMgold}{HTML}{D2B16A}
\definecolor{UNAMteal}{HTML}{16858F}
\definecolor{Pale}{HTML}{EDF3F8}
\definecolor{Muted}{HTML}{496079}
\definecolor{Good}{HTML}{16765A}
\definecolor{Warning}{HTML}{B66B17}

\usetheme{Madrid}
\usecolortheme{default}
\setbeamercolor{structure}{fg=UNAMnavy}
\setbeamercolor{frametitle}{bg=UNAMnavy,fg=white}
\setbeamercolor{title}{fg=white}
\setbeamercolor{author in head/foot}{bg=UNAMnavy,fg=white}
\setbeamercolor{date in head/foot}{bg=UNAMnavy,fg=white}
\setbeamercolor{normal text}{fg=UNAMdeep,bg=white}
\setbeamercolor{block title}{bg=UNAMnavy,fg=white}
\setbeamercolor{block body}{bg=Pale,fg=UNAMdeep}
\setbeamercolor{alerted text}{fg=Warning}
\setbeamertemplate{navigation symbols}{}
\setbeamertemplate{footline}{%
  \leavevmode\hbox{\begin{beamercolorbox}[wd=.84\paperwidth,ht=2.5ex,dp=1ex,leftskip=1.2em]{author in head/foot}%
    Sistemas Inteligentes $\vert$ NYC $\vert$ Fourier y Haar
  \end{beamercolorbox}\begin{beamercolorbox}[wd=.16\paperwidth,ht=2.5ex,dp=1ex,center]{date in head/foot}%
    \insertframenumber / \inserttotalframenumber
  \end{beamercolorbox}}}
\setbeamersize{text margin left=9mm,text margin right=9mm}
\setlength{\parskip}{2pt}
\setlength{\tabcolsep}{4pt}
\newcommand{\clav}[1]{\textcolor{UNAMteal}{\textbf{#1}}}
\newcommand{\porq}[1]{\vspace{1.2mm}\begin{block}{¿Por qué se hizo así?}\small #1\end{block}}
\newcommand{\grafico}[2]{\IfFileExists{figuras/#1}{\includegraphics[width=\linewidth,height=#2,keepaspectratio]{figuras/#1}}{\fbox{\parbox{.80\linewidth}{\centering Imagen opcional: figuras/#1}}}}
\newcommand{\dato}[2]{\begin{block}{#1}\centering\Large\textbf{#2}\end{block}}
\lstset{basicstyle=\ttfamily\footnotesize,backgroundcolor=\color{Pale},frame=single,rulecolor=\color{Pale},showstringspaces=false,breaklines=true,columns=fullflexible}
\hypersetup{colorlinks=true,urlcolor=UNAMteal,linkcolor=UNAMnavy}

\title[NYC: Fourier y Haar]{Predicción de la duración de viajes\\de taxi en Nueva York}
\subtitle{Limpieza, selección de variables, Fourier, Haar y evaluación sin fuga}
\author{Christian Gabriel Villafuerte Aguilar}
\institute{Sistemas Inteligentes \\ Licenciatura en Matemáticas Aplicadas y Computación \\ FES Acatlán -- UNAM}
\date{Octubre de 2026}

\begin{document}

\begin{frame}[plain]
\begin{tikzpicture}[remember picture,overlay]
  \fill[UNAMnavy] (current page.south west) rectangle (current page.north east);
  \fill[UNAMgold] ([yshift=-8mm]current page.north west) rectangle ([yshift=-10mm]current page.north east);
  \node[align=center,text=white,font=\LARGE\bfseries,text width=.90\paperwidth]
      at ([yshift=12mm]current page.center)
      {Predicción de la duración\\de viajes de taxi en Nueva York};
  \node[align=center,text=UNAMgold,font=\large,text width=.90\paperwidth]
      at ([yshift=-6mm]current page.center)
      {Desde la limpieza hasta Fourier, wavelets Haar y ensambles};
  \draw[UNAMteal,line width=1.8pt]
      ([xshift=19mm,yshift=24mm]current page.south west)--([xshift=-19mm,yshift=24mm]current page.south east);
  \node[align=center,text=white,font=\normalsize]
      at ([yshift=17mm]current page.south){Christian Gabriel Villafuerte Aguilar};
  \node[align=center,text=white,font=\small]
      at ([yshift=10mm]current page.south){Sistemas Inteligentes $\cdot$ MAC $\cdot$ FES Acatlán, UNAM};
  \node[align=center,text=white,font=\small]
      at ([yshift=4mm]current page.south){Octubre de 2026};
\end{tikzpicture}
\end{frame}

\section{Comprensión del problema}
\begin{frame}{Ruta de la exposición}
\begin{enumerate}
  \item Qué se quiere predecir y cuál es la métrica oficial.
  \item Auditoría de datos, limpieza y partición correcta.
  \item Construcción de \clav{81 variables} e imputación/escalado.
  \item Comparación y ajuste de modelos de regresión.
  \item Evaluación y optimización geográfica, Fourier y wavelets Haar.
  \item Ensamble, holdout y entrega del CSV.
\end{enumerate}
\porq{El objetivo no es mostrar solamente un puntaje: se debe justificar un proceso reproducible y evitar fuga de información.}
\end{frame}

\begin{frame}{¿Cuál es el problema?}
\begin{columns}[T]
  \column{.57\textwidth}
    \begin{itemize}
      \item Entrada: datos de cada viaje de taxi amarillo.
      \item Salida deseada: \texttt{trip\_duration}, en \textbf{segundos}.
      \item Entrenamiento: \textbf{700,000} viajes con duración conocida.
      \item Evaluación ciega: \textbf{100,000} viajes sin duración.
      \item Entrega: un valor no negativo por cada \texttt{id}.
    \end{itemize}
  \column{.39\textwidth}
    \dato{Tipo de aprendizaje}{Supervisado}
    \dato{Tipo de tarea}{Regresión}
\end{columns}
\porq{La duración puede tomar muchos valores numéricos; no estamos decidiendo una clase ``corto/largo''. Por eso no corresponde usar un clasificador.}
\end{frame}

\begin{frame}{¿Por qué RMSLE y no accuracy o F1?}
\begin{block}{Métrica oficial: Root Mean Squared Logarithmic Error}
\[
\mathrm{RMSLE}=\sqrt{\frac1n\sum_{i=1}^n\Big[\ln(1+\widehat y_i)-\ln(1+y_i)\Big]^2}.
\]
\end{block}
\begin{columns}[T]
\column{.48\textwidth}
  \textbf{En regresión:} RMSLE, MAE, RMSE, $R^2$, MAPE.
\column{.48\textwidth}
  \textbf{En clasificación:} accuracy, precision, recall, F1.
\end{columns}
\vspace{2mm}
\begin{itemize}
  \item \clav{Menor RMSLE significa mejor resultado.}
  \item La escala logarítmica da relevancia a los errores relativos.
\end{itemize}
\porq{Crear clases artificiales para calcular F1 cambiaría el problema y no optimizaría la evaluación del concurso.}
\end{frame}

\begin{frame}{Metodología de clase: ciclo de vida CRISP-DM}
\centering
\begin{tikzpicture}[>=Latex,node distance=3mm,every node/.style={font=\small}]
\tikzset{etapa/.style={draw=UNAMnavy,thick,rounded corners=3pt,fill=Pale,minimum height=11mm,text width=22mm,align=center}}
\node[etapa] (a) {Entender el\\ problema};
\node[etapa,right=of a] (b) {Auditar\\ datos};
\node[etapa,right=of b] (c) {Preparar\\ variables};
\node[etapa,right=of c] (d) {Entrenar\\ modelos};
\node[etapa,right=of d] (e) {Evaluar\\ resultados};
\draw[->,very thick,UNAMteal] (a)--(b);
\draw[->,very thick,UNAMteal] (b)--(c);
\draw[->,very thick,UNAMteal] (c)--(d);
\draw[->,very thick,UNAMteal] (d)--(e);
\draw[->,UNAMgold,thick] (e.south) .. controls +(0,-13mm) and +(0,-13mm) .. node[below,font=\scriptsize]{revisar variables y repetir si hace falta} (c.south);
\end{tikzpicture}
\vspace{8mm}
\begin{block}{Aplicación concreta en el notebook}
CSV $\rightarrow$ auditoría $\rightarrow$ partición $\rightarrow$ 81 atributos $\rightarrow$ modelos $\rightarrow$ holdout $\rightarrow$ \texttt{submission.csv}.
\end{block}
\porq{Es el enfoque explicado en la presentación \emph{Ciclo de vida del ML}: antes del algoritmo van la calidad de datos y el planteamiento de la métrica.}
\end{frame}

\section{Datos y limpieza}
\begin{frame}{¿Qué información contienen los archivos?}
\small
\begin{tabular}{p{.28\linewidth}p{.66\linewidth}}
\toprule
\textbf{Campo} & \textbf{Interpretación y uso}\tabularnewline\midrule
\texttt{id} & Identificador: se conserva para entregar, no se entrena con él.\\
\texttt{vendor\_id} & Empresa: se mantiene como variable numérica/código.\\
\texttt{pickup\_datetime} & Fecha y hora de inicio: genera variables temporales.\\
\texttt{passenger\_count} & Número declarado de pasajeros: contiene alguna anomalía.\\
\texttt{pickup\_*}, \texttt{dropoff\_*} & Coordenadas de origen/destino: permiten calcular distancias.\\
\texttt{store\_and\_fwd\_flag} & Conexión interrumpida (Y/N): se codifica 1/0.\\
\texttt{trip\_duration} & Variable objetivo, solo en \texttt{train.csv}.\\
\bottomrule
\end{tabular}
\porq{No se usa \texttt{id} como predictor porque es una etiqueta administrativa sin una relación física justificada con la duración.}
\end{frame}

\begin{frame}{Auditoría inicial: ¿los CSV están limpios?}
\begin{columns}[T]
\column{.5\textwidth}
  \begin{itemize}
    \item \textbf{0} valores nulos originales en \texttt{train} y \texttt{test}.
    \item \textbf{0} IDs duplicados en ambos archivos.
    \item \textbf{0} duraciones menores o iguales a cero.
    \item \textbf{24} registros con cero pasajeros.
  \end{itemize}
\column{.46\textwidth}
  \begin{itemize}
    \item \textbf{984} duraciones mayores que 12 horas.
    \item \textbf{2} duraciones mayores que 24 horas.
    \item \textbf{32} orígenes fuera de caja GPS NYC.
    \item \textbf{54} destinos fuera de esa caja.
  \end{itemize}
\end{columns}
\porq{Que no haya nulos no significa que todos los registros sean plausibles. También se revisan consistencia temporal, GPS y valores extremos.}
\end{frame}

\begin{frame}{Exploración de datos: antes de modelar}
\centering\grafico{01_exploracion.png}{.66\textheight}
\vspace{1mm}
{\footnotesize Distribución logarítmica de duración, número de viajes por hora y ubicaciones de origen.}
\porq{La duración tiene una cola larga; una gráfica en segundos ocultaría gran parte de los viajes usuales. Las horas y ubicaciones justifican crear variables de contexto.}
\end{frame}

\begin{frame}{Decisiones explícitas sobre anomalías}
\small
\begin{tabular}{p{.25\linewidth}p{.40\linewidth}p{.28\linewidth}}
\toprule
\textbf{Caso} & \textbf{Acción aplicada} & \textbf{Motivo}\tabularnewline\midrule
Viaje de más de 12 h & Se excluye \textbf{solo del ajuste}; no de validación ni holdout. & Reducir influencia de etiquetas extremas sin maquillar la evaluación.\\
GPS fuera de zona & Se marca con indicador y se convierte su coordenada a NaN. & No calcular distancias absurdas.\\
Cero pasajeros & Se marca con bandera de calidad; se conserva el dato. & Señal de calidad sin borrar el viaje.\\
\texttt{test.csv} & \textbf{Nunca} se eliminan filas. & Deben existir 100,000 predicciones.\\
\bottomrule
\end{tabular}
\porq{El umbral de 12 horas es una decisión heurística del experimento; conviene validarlo en estudios posteriores, no afirmar que es una regla universal.}
\end{frame}

\begin{frame}{División correcta: entrenamiento, validación y holdout}
\centering
\begin{tikzpicture}
\fill[UNAMnavy] (0,0) rectangle (10.3,1.15);
\fill[UNAMteal] (10.3,0) rectangle (11.6,1.15);
\fill[UNAMgold] (11.6,0) rectangle (12.9,1.15);
\node[white,font=\normalsize] at (5.1,.55){80\% $\rightarrow$ 560,000 viajes};
\node[white,font=\scriptsize,align=center] at (10.95,.55){10\%\\70,000};
\node[UNAMdeep,font=\scriptsize,align=center] at (12.25,.55){10\%\\70,000};
\node[below] at (5.1,-.12){Entrenamiento};
% Ambos grupos están identificados dentro de los bloques.
\end{tikzpicture}
\vspace{8mm}
\begin{itemize}
\item En los 560,000 de la partición de entrenamiento se filtraron 788 viajes $>12$ h: \clav{559,212 usados para ajustar}.
\item Validación: compara modelos y parámetros.
\item Holdout: se consulta una vez, tras la selección.
\end{itemize}
\porq{Si el holdout participa en el ajuste, deja de ser una evaluación independiente. Semilla fija: \texttt{random\_state=42}.}
\end{frame}

\section{Preparación y variables}
\begin{frame}{Del dato crudo a una matriz numérica}
\centering
\begin{tikzpicture}[>=Latex,node distance=5mm,every node/.style={font=\scriptsize}]
\tikzset{paso/.style={draw=UNAMnavy,fill=Pale,rounded corners,minimum width=22mm,minimum height=14mm,align=center}}
\node[paso] (a) {CSV\\ original};
\node[paso,right=of a] (b) {Fecha y GPS\\ consistentes};
\node[paso,right=of b] (c) {81\ atributos\\ derivados};
\node[paso,right=of c] (d) {Imputar\\ medianas};
\node[paso,right=of d] (e) {Matriz\\ $X$};
\foreach \x/\y in {a/b,b/c,c/d,d/e}{\draw[->,thick,UNAMteal] (\x)--(\y);}
\end{tikzpicture}
\vspace{8mm}
\begin{block}{Propiedad importante}
La función \texttt{create\_features} usa únicamente datos disponibles en los archivos de origen, \textbf{nunca la duración real} para construir predictores.
\end{block}
\porq{El modelo no comprende de forma automática que dos coordenadas implican una distancia o que las 23:00 y las 00:00 están muy próximas.}
\end{frame}

\begin{frame}{Las 81 características: ¿de dónde salen?}
\footnotesize
\begin{tabular}{p{.25\linewidth}p{.45\linewidth}p{.22\linewidth}}
\toprule
\textbf{Grupo} & \textbf{Ejemplos implementados} & \textbf{Qué aportan}\tabularnewline\midrule
Tiempo & Hora, minuto, día, fin de semana, hora pico, seno/coseno & Congestión y periodicidad\\
Geometría & Haversine, Manhattan aprox., $\Delta x$, $\Delta y$, rumbo & Longitud y dirección\\
Zonas & Rejillas GPS, puntos medios & Barrios y trayectos frecuentes\\
Referencias & Distancia a JFK, LGA, EWR, Times Square y otros & Contexto urbano\\
Interacciones & Hora $\times$ distancia, hora pico $\times$ distancia & Relaciones no lineales\\
Calidad & GPS inválido, pasajero inválido, conexión diferida & Anomalías del registro\\
\bottomrule
\end{tabular}
\porq{Estas 81 columnas resultan de combinar conocimiento del dominio con transformaciones reproducibles; su cantidad no garantiza por sí sola calidad.}
\end{frame}

\begin{frame}{Variables temporales: ¿por qué usar seno y coseno?}
\begin{columns}[T]
\column{.52\textwidth}
\[
 h_{\sin}=\sin\left(\frac{2\pi h}{24}\right),\qquad
 h_{\cos}=\cos\left(\frac{2\pi h}{24}\right).
\]
\begin{itemize}
\item 23:00 y 00:00 son horas contiguas.
\item Una hora numérica normal no refleja bien esa cercanía.
\item También se aplican ciclos a días y meses.
\end{itemize}
\column{.44\textwidth}
\centering
\begin{tikzpicture}[scale=1.15]
\draw[->,gray] (-1.35,0)--(1.35,0);
\draw[->,gray] (0,-1.35)--(0,1.35);
\draw[thick,UNAMnavy] (0,0) circle (1);
\fill[UNAMteal] (1,0) circle (.07);
\fill[UNAMgold] ({cos(-15)},{sin(-15)}) circle (.075);
\node[font=\scriptsize,right] at (1,.12){00 h};
\node[font=\scriptsize,right] at (0.84,-.37){23 h};
\end{tikzpicture}
\end{columns}
\porq{La codificación cíclica hace que el cambio de día no aparezca como un salto artificial de 23 a 0.}
\end{frame}

\begin{frame}{¿Cómo calcular la distancia entre dos coordenadas?}
\begin{block}{Fórmula de Haversine (distancia sobre la esfera terrestre)}
\small
\[
\begin{aligned}
 a&=\sin^2\!\left(\frac{\Delta\varphi}{2}\right)
 +\cos\varphi_1\cos\varphi_2\sin^2\!\left(\frac{\Delta\lambda}{2}\right),\\[-1mm]
 d&=2R\arcsin(\sqrt{a}), \qquad R\simeq 6371\,\mathrm{km}.
\end{aligned}
\]
\end{block}
\begin{itemize}
\item $\varphi$: latitud; $\lambda$: longitud, en radianes.
\item La distancia Haversine estima el trayecto directo, \textbf{no} las calles recorridas.
\item Se añade distancia Manhattan aproximada para representar desplazamientos ortogonales.
\end{itemize}
\porq{Las coordenadas originales son difíciles de interpretar solas; las distancias aportan una relación más directa con el tiempo esperado.}
\end{frame}

\begin{frame}{Geografía urbana y relaciones entre variables}
\begin{columns}[T]
\column{.49\textwidth}
\begin{block}{Origen / destino}
\begin{itemize}
\item Distancia a JFK, LaGuardia y Newark.
\item Distancia a Times Square, Grand Central y Downtown.
\item Celdas y punto medio de la trayectoria.
\end{itemize}
\end{block}
\column{.49\textwidth}
\begin{block}{Combinaciones}
\begin{itemize}
\item \texttt{hour\_x\_distance}.
\item \texttt{rush\_x\_distance}.
\item Ángulo del trayecto y sus seno/coseno.
\item Distancias al cuadrado y logarítmicas.
\end{itemize}
\end{block}
\end{columns}
\porq{Un trayecto de 5 km durante hora pico puede durar distinto que el mismo trayecto por la madrugada; la interacción captura esa dependencia.}
\end{frame}

\begin{frame}{Imputación: ¿qué hacer cuando una coordenada es inválida?}
\begin{enumerate}
\item Detectar coordenadas fuera de una caja geográfica razonable.
\item Crear bandera de calidad y sustituir la coordenada inválida por \texttt{NaN}.
\item Aprender la \textbf{mediana de cada característica} usando \textbf{únicamente entrenamiento}.
\item Aplicar las mismas medianas a validación, holdout y test.
\end{enumerate}
\begin{block}{Implementación real}
\texttt{SimpleImputer(strategy='median', keep\_empty\_features=True)}
\end{block}
\porq{La mediana suele ser más resistente que la media a extremos. Si se calculase usando todo el conjunto, entrarían estadísticas de validación o prueba en el entrenamiento.}
\end{frame}

\begin{frame}{Normalización, escalado y logaritmo: NO son lo mismo}
\small
\begin{tabular}{p{.28\linewidth}p{.32\linewidth}p{.30\linewidth}}
\toprule
\textbf{Técnica} & \textbf{Expresión} & \textbf{¿Se usó?}\tabularnewline\midrule
Estandarización & $(x-\mu)/\sigma$ & Se explica, pero no se aplicó a los árboles.\\
Min--max & $(x-x_{\min})/(x_{\max}-x_{\min})$ & No se aplicó.\\
Escalado robusto & $(x-\mathrm{mediana})/\mathrm{IQR}$ & \textbf{Sí: solo en Ridge}.\\
Transformación del objetivo & $z=\ln(1+y)$ & \textbf{Sí: en todos los modelos}.\\
\bottomrule
\end{tabular}
\vspace{2mm}
\porq{Ridge penaliza coeficientes y necesita escalas comparables; los árboles separan por umbrales, así que no requieren escalado para funcionar.}
\end{frame}

\begin{frame}{Transformación de la duración para entrenar}
\begin{columns}[T]
\column{.52\textwidth}
\begin{block}{Objetivo de entrenamiento}
\[
 z=\ln(1+\texttt{trip\_duration}).
\]
\end{block}
\begin{block}{Recuperar segundos}
\[
 \widehat y=\max\{0,\exp(\widehat z)-1\}.
\]
\end{block}
\column{.44\textwidth}
\begin{itemize}
\item Reduce la asimetría de la duración.
\item Es consistente con la métrica RMSLE.
\item Evita valores negativos en el archivo final.
\item \textbf{No} significa que se entregue el logaritmo.
\end{itemize}
\end{columns}
\porq{Se optimiza en la escala logarítmica y después se regresa a segundos, que es la unidad requerida por el reto.}
\end{frame}

\begin{frame}{¿Cómo se evitó la fuga de información?}
\centering
\begin{tikzpicture}[node distance=7mm,>=Latex,every node/.style={font=\small}]
\node[draw=UNAMnavy,thick,fill=Pale,rounded corners,minimum width=41mm,minimum height=10mm] (tr) {Entrenamiento \texttt{tr\_fit}};
\node[draw=UNAMteal,thick,fill=Pale,rounded corners,right=17mm of tr,minimum width=53mm,minimum height=10mm] (imp) {\texttt{imputer.fit(Xtr)}};
\draw[->,thick,UNAMteal] (tr)--(imp);
\node[draw=UNAMnavy,fill=white,below=5mm of tr,minimum width=41mm,minimum height=8mm] (va) {Validación / holdout / test};
\node[draw=UNAMteal,fill=white,below=5mm of imp,minimum width=53mm,minimum height=8mm] (tx) {\texttt{imputer.transform(X)}};
\draw[->,thick,UNAMteal] (va)--(tx);
\draw[->,thick,UNAMgold] (imp)--(tx);
\end{tikzpicture}
\vspace{8mm}
\begin{itemize}
\item No se ajustan medianas ni \texttt{RobustScaler} con validación o test.
\item No se calcula velocidad usando duración real ni se inventa hora de llegada.
\item La evaluación holdout ocurre después de seleccionar el ganador.
\end{itemize}
\porq{Una métrica artificialmente baja por fuga de información no mide la capacidad real de generalizar.}
\end{frame}

\section{Modelado y selección}
\begin{frame}{Primero: dos referencias fáciles de explicar}
\begin{columns}[T]
\column{.48\textwidth}
\begin{block}{Modelo ingenuo: mediana}
Predice para todos los viajes la mediana de la duración en escala logarítmica.
\vspace{2mm}
\centering\Large\textbf{RMSLE = 0.795413}
\end{block}
\column{.48\textwidth}
\begin{block}{Ridge: regresión lineal}
Modelo lineal con regularización $L_2$ y \texttt{RobustScaler}.
\vspace{2mm}
\centering\Large\textbf{RMSLE = 0.462847}
\end{block}
\end{columns}
\porq{El baseline fija el mínimo que debemos superar; Ridge evalúa cuánto se puede explicar con una relación lineal regularizada.}
\end{frame}

\begin{frame}{¿Por qué Ridge en vez de quedarse con regresión lineal?}
\begin{block}{Idea matemática}
\[
\widehat{\boldsymbol\beta}
 =\arg\min_{\boldsymbol\beta}\left\{\sum_{i=1}^n(z_i-\beta_0-\mathbf x_i^\top\boldsymbol\beta)^2
 +\alpha\lVert\boldsymbol\beta\rVert_2^2\right\},\qquad \alpha=100.
\]
\end{block}
\begin{itemize}
\item La penalización reduce coeficientes excesivamente grandes.
\item Es útil con muchas variables correlacionadas.
\item Se ajusta sobre atributos escalados de forma robusta.
\item Pero tiene dificultades con interacciones no lineales complejas.
\end{itemize}
\porq{Ridge es una referencia interpretable y regularizada, no necesariamente el modelo final.}
\end{frame}

\begin{frame}{¿Por qué pasar a árboles de boosting?}
\begin{columns}[T]
\column{.51\textwidth}
\begin{itemize}
\item Distancia y duración no mantienen siempre relación lineal.
\item Influyen hora, zona, aeropuerto y trayecto simultáneamente.
\item \textbf{Boosting} añade árboles que corrigen errores anteriores.
\item Puede capturar interacciones y cortes no lineales.
\end{itemize}
\column{.45\textwidth}
\centering
\begin{tikzpicture}[scale=.87,>=Latex]
\node[draw=UNAMnavy,fill=Pale,rounded corners] (a) at (0,1.3){Árbol 1};
\node[draw=UNAMnavy,fill=Pale,rounded corners] (b) at (2.2,1.3){Árbol 2};
\node[draw=UNAMnavy,fill=Pale,rounded corners] (c) at (4.4,1.3){Árbol 3};
\draw[->,thick,UNAMteal] (a)--(b);
\draw[->,thick,UNAMteal] (b)--(c);
\draw[->,thick,UNAMgold] (2.2,.8)--(2.2,0) node[below,font=\small,align=center]{suma gradual\\de correcciones};
\end{tikzpicture}
\end{columns}
\porq{Los métodos de árboles trabajan con umbrales; por ello aquí no añadimos escalado de entrada que no aporta una ventaja clara.}
\end{frame}

\begin{frame}{Modelos probados y motivo de comparación}
\footnotesize
\begin{tabular}{p{.23\linewidth}p{.43\linewidth}p{.27\linewidth}}
\toprule
\textbf{Modelo} & \textbf{Idea} & \textbf{Motivo del ensayo}\tabularnewline\midrule
LightGBM & Boosting eficiente con árboles orientados a hojas. & Eficiencia en datos tabulares grandes.\\
XGBoost & Boosting regularizado con árboles. & Comparación con otra implementación sólida.\\
CatBoost & Boosting con mecanismos de regularización propios. & Contrastar estabilidad y desempeño.\\
\midrule
Blends & Combinaciones ponderadas en escala log. & Comprobar si se complementan sus errores.\\
\bottomrule
\end{tabular}
\porq{No se declaró ganador por popularidad: todos se compararon con el mismo conjunto de validación y la misma métrica.}
\end{frame}

\begin{frame}{¿Qué parámetros se ajustaron en LightGBM?}
\scriptsize
\begin{tabular}{p{.36\linewidth}rrr}
\toprule
\textbf{Hiperparámetro} & \textbf{Equilibrado} & \textbf{Profundo} & \textbf{Regularizado}\\\midrule
\texttt{n\_estimators} & 1400 & 1600 & 1650\\
\texttt{learning\_rate} & 0.045 & 0.045 & 0.035\\
\texttt{num\_leaves} & 63 & 105 & 39\\
\texttt{min\_child\_samples} & 45 & 75 & 75\\
\texttt{colsample\_bytree} & 0.87 & 0.92 & 0.83\\
\texttt{subsample} & 0.87 & 0.90 & 0.88\\
\texttt{reg\_lambda} & 2.5 & 5 & 6\\\midrule
\textbf{RMSLE de validación} & \textbf{0.379249} & \textbf{0.375735} & \textbf{0.382771}\\
\bottomrule
\end{tabular}
\vspace{2mm}
\begin{block}{Criterio aplicado}
Entrenamiento sobre $\ln(1+y)$; evaluación por RMSLE; \emph{early stopping} con 90 rondas sin mejora para LightGBM.
\end{block}
\porq{Más hojas pueden capturar estructuras complejas, pero se controlan con mínimo de muestras por hoja y regularización. No se exploró todo el espacio de parámetros.}
\end{frame}

\begin{frame}{Resultados de validación: comparación objetiva}
\begin{columns}[T]
\column{.53\textwidth}
\scriptsize
\begin{tabular}{lr}
\toprule Modelo & RMSLE $\downarrow$\\\midrule
\textbf{LightGBM profundo} & \textbf{0.375735}\\
LightGBM equilibrado & 0.379249\\
XGBoost & 0.381371\\
LightGBM regularizado & 0.382771\\
CatBoost & 0.385040\\
Ridge escalado & 0.462847\\
Mediana & 0.795413\\
\bottomrule
\end{tabular}
\par\vspace{3mm}
\noindent\textbf{Ganador entre los modelos ensayados:} LightGBM profundo.
\column{.43\textwidth}
\begin{tikzpicture}
\begin{axis}[xbar,xmin=0,xmax=.85,width=\linewidth,height=54mm,bar width=7pt,axis x line=bottom,axis y line=none,ytick=data,yticklabels={Mediana,Ridge,CatBoost,XGBoost,LightGBM},tick label style={font=\scriptsize},nodes near coords,point meta=explicit symbolic,every node near coord/.append style={font=\tiny,anchor=west},enlarge y limits=.14]
\addplot[fill=UNAMteal,draw=UNAMteal] coordinates {(.7954,0) [.795] (.4628,1) [.463] (.3850,2) [.385] (.3814,3) [.381] (.3757,4) [.376]};
\end{axis}
\end{tikzpicture}
\end{columns}
\porq{El único criterio de selección fue minimizar RMSLE en la validación; valores menores indican mejores predicciones relativas.}
\end{frame}

\begin{frame}{¿Se intentó combinar modelos? Sí}
\small
\begin{tabular}{p{.75\linewidth}r}
\toprule
\textbf{Ensamble probado} & \textbf{RMSLE}\\\midrule
75\% LightGBM profundo + 25\% LightGBM equilibrado & 0.375912\\
50\% LightGBM profundo + 50\% LightGBM equilibrado & 0.376559\\
25\% LightGBM profundo + 75\% LightGBM equilibrado & 0.377673\\
50/30/20 (profundo / equilibrado / XGBoost) & 0.376475\\
\midrule
\textbf{LightGBM profundo sin mezcla} & \textbf{0.375735}\\\bottomrule
\end{tabular}
\vspace{2mm}
\porq{Un ensamble no garantiza mejora. Como ninguno superó al modelo individual, el resultado final usa un solo LightGBM: mejor RMSLE y menor complejidad.}
\end{frame}

\section{Evaluación de la línea base}
\begin{frame}{Primera evaluación: 70,000 viajes holdout}
\begin{columns}[T]
\column{.49\textwidth}
\dato{RMSLE en validación}{0.375735}
\column{.49\textwidth}
\dato{RMSLE en holdout}{0.381219}
\end{columns}
\vspace{2mm}
\begin{center}
\small
\begin{tabular}{ll}
\toprule\textbf{Métrica complementaria en holdout} & \textbf{Valor}\\\midrule
MAE & 270.92 segundos\\
RMSE & 3,062.55 segundos\\
$R^2$ & 0.0390\\
MAPE & 25.48\%\\\bottomrule
\end{tabular}
\end{center}
\porq{El holdout se mantuvo ajeno a la selección, por lo que su RMSLE es una estimación más honesta de generalización; no es el puntaje de test ciego.}
\end{frame}

\begin{frame}{Interpretación: RMSLE vs. MAE, RMSE y $R^2$}
\small
\begin{tabular}{p{.19\linewidth}p{.71\linewidth}}
\toprule
\textbf{Métrica} & \textbf{Qué nos dice}\tabularnewline\midrule
RMSLE & Penaliza errores al comparar \textbf{log-duraciones}; es la métrica oficial.\\
MAE & En promedio hay \textbf{270.92 s} de error absoluto en holdout (unos 4.5 min).\\
RMSE & \textbf{3,062.55 s}; muy sensible a errores grandes en segundos.\\
$R^2$ & \textbf{0.039}: el ajuste en escala de segundos explica poca varianza de estas etiquetas.\\
MAPE & \textbf{25.48\%}, pero puede ser inestable cuando la duración real es pequeña.\\
\bottomrule
\end{tabular}
\porq{El mismo modelo puede rendir bien en RMSLE y mal en RMSE/$R^2$: no son métricas intercambiables y los valores extremos influyen mucho en segundos.}
\end{frame}

\begin{frame}{Diagnóstico: ¿en qué casos se equivoca?}
\centering\grafico{03_diagnostico_holdout.png}{.55\textheight}
\vspace{1mm}
{\footnotesize Predicción frente a valor real en logaritmos; distribución del error logarítmico.}
\porq{Las dispersiones y sesgos ayudan a detectar sobreestimación, subestimación y viajes atípicos. Una sola métrica agregada no explica todos los fallos.}
\end{frame}

\begin{frame}{Importancia de variables según LightGBM}
\begin{columns}[T]
\column{.54\textwidth}
\centering\grafico{04_importancia_variables.png}{.61\textheight}
\column{.43\textwidth}
\small
\begin{itemize}
\item Distancia a \textbf{JFK} desde el origen.
\item Componentes cíclicas de la \textbf{hora}.
\item \textbf{Día del año}.
\item Distancia a JFK desde destino.
\item \textbf{Hora $\times$ distancia}.
\item \textbf{Distancia Haversine}.
\end{itemize}
\end{columns}
\porq{Es una importancia interna del árbol (cantidad de uso de variables en cortes), no demuestra causalidad ni ordena necesariamente el efecto real de cada característica.}
\end{frame}

\begin{frame}{El error depende de la distancia del trayecto}
\small
\begin{tabular}{lrr}
\toprule\textbf{Distancia Haversine} & \textbf{Viajes holdout} & \textbf{RMSLE}\\\midrule
Menos de 0.5 km & 2,504 & 1.103559\\
0.5 a 2 km & 31,070 & 0.368468\\
2 a 5 km & 23,896 & 0.290010\\
5 a 10 km & 8,250 & 0.260631\\
10 a 25 km & 4,200 & 0.275809\\
Más de 25 km & 76 & 0.245508\\\bottomrule
\end{tabular}
\vspace{1mm}
{\scriptsize Nota: cuatro viajes holdout no entran en estos rangos por GPS inválido.}\par
\vspace{1mm}
\porq{En viajes muy cortos, paradas, espera y registros de GPS pueden dominar el tiempo; algunos grupos tienen pocas observaciones. Se trata de un diagnóstico, no de una prueba causal.}
\end{frame}


% --- SEGUNDA ETAPA: GEOGRAFIA, FOURIER Y HAAR ---
\section{Optimización geográfica sin fuga}
\begin{frame}{Segunda etapa: ¿por qué no bastaron las 81 variables?}
\begin{columns}[T]
\column{.53\textwidth}
\small
\begin{itemize}
\item La distancia recta no describe el trazado urbano ni la congestión.
\item Se probaron agrupaciones geográficas, ejes de Manhattan y estadísticas de trayectos.
\item Algunas codificaciones de zonas \alert{empeoraron} el RMSLE: eran demasiado específicas.
\item Se conservaron variables geométricas continuas y estadísticas \emph{out-of-fold} (OOF).
\end{itemize}
\column{.44\textwidth}
\begin{block}{Criterio de decisión}\centering
\textbf{Mejora en validación}\\[1mm]
NO complejidad por sí sola
\end{block}
\begin{block}{Experimento previo}
\centering \large 0.374624
\par\small RMSLE de validación\end{block}
\end{columns}
\porq{Probamos hipótesis de dominio y retiramos características que agregaban sobreajuste. La evaluación usa siempre la misma partición.}
\end{frame}

\begin{frame}{Geografía urbana: ejes, regiones y trayectos}
\begin{columns}[T]
\column{.56\textwidth}
\small
\begin{itemize}
\item Conversión aproximada de latitud/longitud a desplazamientos en kilómetros.
\item Proyección rotada respecto de los ejes principales de Manhattan.
\item Diferencias de coordenadas entre origen y destino.
\item Clústeres y parejas origen--destino, pero con regularización/codificación OOF.
\end{itemize}
\column{.40\textwidth}
\begin{block}{Ejemplo geométrico}
\[ \Delta x = 111.195\cos(\bar\varphi)\Delta\lambda\]
\[ \Delta y = 111.195\Delta\varphi \]
\[ d_{\mathrm{man}}\approx |\Delta x|+|\Delta y|\]
\end{block}
\end{columns}
\porq{La red vial hace que dos trayectos con similar Haversine puedan tener tiempos distintos. Las nuevas coordenadas no emplean el objetivo.}
\end{frame}

\begin{frame}{Codificación de promedios geográficos \emph{out-of-fold}}
\small
Para usar tiempos históricos por zona, cada viaje debe recibir una estimación \textbf{sin su propia etiqueta}. 
\vspace{2mm}
\begin{center}
\begin{tikzpicture}[font=\scriptsize,>=Latex]
\foreach \i in {0,...,4}{\draw[fill=UNAMteal!16,draw=UNAMteal] (1.65*\i,1.12) rectangle +(1.48,.49);
\node at (1.65*\i+.74,1.37) {Fold \the\numexpr\i+1\relax};}
\draw[->,very thick,UNAMnavy] (4.0,1.0)--(4.0,.63);
\node[draw=UNAMnavy,rounded corners,fill=Pale,align=center] at (4,.27) {Ajustar con 4 folds $\to$ transformar el fold excluido};
\end{tikzpicture}
\end{center}
\[
\widehat\mu_g=\frac{n_g\overline z_g+\alpha\overline z}{n_g+\alpha},\qquad z=\log(1+y),
\]
\begin{itemize}
\item \(\alpha\) suaviza categorías con muy pocos registros.
\item Se repite el procedimiento para los cinco folds.
\item Validación/test usan mapeos entrenados \textbf{solo} con train.
\end{itemize}
\end{frame}

\begin{frame}{Resultados de la mejora geográfica}
\small
\begin{center}
\begin{tabular}{lcc}
\toprule
\textbf{Configuración} & \textbf{Validación RMSLE} & \textbf{Holdout RMSLE}\\
\midrule
LightGBM inicial profundo & 0.375735 & 0.381219\\
Ensamble geográfico + OOF & \textbf{0.374624} & \textbf{0.379288}\\
\bottomrule
\end{tabular}
\end{center}
\vspace{2mm}
\begin{block}{Lección}
El avance existe, pero es modesto: aún hay estructura temporal y variaciones locales que las coordenadas por sí solas no explican.
\end{block}
\porq{Por eso probamos análisis espectral y multirresolución como una \emph{nueva hipótesis} y no como una garantía de RMSLE 0.20.}
\end{frame}

\section{Análisis espectral de Fourier}
\begin{frame}{¿Por qué Fourier para viajes de taxi?}
\small
\begin{itemize}
\item La demanda cambia con \textbf{horas pico, turnos, fines de semana} y estacionalidad.
\item Una hora numérica no codifica automáticamente que las 23:59 y 00:01 están muy cerca.
\item Fourier describe patrones periódicos mediante senos y cosenos a distintas frecuencias.
\end{itemize}
\[
\phi_{k,\sin}(t)=\sin\!\left(\frac{2\pi k t}{T}\right),\quad
\phi_{k,\cos}(t)=\cos\!\left(\frac{2\pi k t}{T}\right).
\]
\begin{block}{Periodos empleados} \(T=24\) horas, \(T=7\) días y una fase estacional sobre 366 días; \(k\) representa la armónica.\end{block}
\end{frame}

\begin{frame}{FFT: pasar de tiempo a frecuencia}
\small
Partimos de una señal de \textbf{conteos de inicios de viaje cada 30 minutos} creada con el subconjunto de entrenamiento.
\[
X_k=\sum_{n=0}^{N-1}x_n e^{-2\pi i kn/N},\qquad k=0,\ldots,N-1.
\]
\begin{itemize}
\item \(x_n\): número transformado de viajes en la franja \(n\).
\item \(|X_k|^2\): energía/potencia asociada a la frecuencia \(k\).
\item FFT: algoritmo eficiente para calcular la transformada discreta.
\end{itemize}
\porq{Se analiza la periodicidad \emph{de los registros de entrenamiento}, sin recurrir a ninguna duración del test ciego.}
\end{frame}

\begin{frame}{Demanda temporal observada en entrenamiento}
\centering
\grafico{senal_demanda_temporal.png}{.68\textheight}
\vspace{1mm}
{\scriptsize Conteos de pickups; las oscilaciones reflejan el ritmo de viajes, no una medición de tráfico real.}
\end{frame}

\begin{frame}{Espectro encontrado: ciclos de 24 h y 7 días}
\centering
\grafico{espectro_fourier.png}{.67\textheight}
\vspace{1mm}
{\scriptsize Se identificaron periodicidades relevantes, incluyendo componentes próximas a 12 h, 24 h y 168 h; \textbf{esto no prueba que mejoren por sí solas la predicción}.}
\end{frame}

\begin{frame}{45 bases deterministas: temporales y espaciales}
\small
\begin{tabular}{p{.32\textwidth}p{.57\textwidth}}
\toprule
\textbf{Grupo} & \textbf{Qué se construyó}\\
\midrule
Fourier diario & Armónicas \(k=2,3,4,6,8,12\): senos y cosenos\\
Fourier semanal & \(k=1,2,3,4,6,8,12\): senos y cosenos\\
Fourier estacional & Armónicas \(k=2,3,4\) del día del año\\
Fases Haar fijas & 5 señales de cambio de signo por bloques\\
Fourier espacial & Cosenos de latitud/longitud de origen y destino a 2 y 6 km\\
\bottomrule
\end{tabular}
\vspace{2mm}
\begin{block}{Ninguna usa la duración}Estas 45 características dependen exclusivamente de \texttt{pickup\_datetime} y de coordenadas disponibles en ambos CSV.\end{block}
\end{frame}

\begin{frame}{Filtrado espectral: tendencia y bandas}
\small
A partir de la serie de demanda, se conservan frecuencias bajas y se reconstruye por transformada inversa:
\[
L_T(t)=\mathcal{F}^{-1}\big(\mathcal{F}(x)\,\mathbf{1}_{\{|f|\le 1/T\}}\big)(t),
\qquad B_T=L_T-L_{T'},
\]
\begin{itemize}
\item Ventanas representadas: \textbf{12, 48 y 168 horas}.
\item El componente lento resume tendencias; la diferencia entre escalas aproxima una banda.
\item Los filtros se aprenden de \textbf{conteos de train}; para aplicar a test se reutiliza esa señal.
\end{itemize}
\porq{Los modelos basados en árboles pueden usar tanto armónicas explícitas como señales suavizadas sin que debamos estandarizarlas.}
\end{frame}

\section{Wavelets Haar y estadísticas históricas}
\begin{frame}{¿Qué agregan las wavelets que Fourier no explica?}
\begin{columns}[T]
\column{.48\textwidth}
\begin{block}{Fourier}\small
Describe patrones \textbf{globales} repetidos; una frecuencia aparece en toda la ventana analizada.
\end{block}
\column{.48\textwidth}
\begin{block}{Haar}\small
Describe \textbf{cambios locales} y diferencias entre resoluciones temporales.
\end{block}
\end{columns}
\vspace{3mm}
\begin{itemize}
\item Un cambio abrupto de demanda puede quedar diluido en un promedio diario.
\item Las aproximaciones muestran la actividad lenta y los detalles muestran transiciones.
\end{itemize}
\porq{Son representaciones \emph{complementarias}; ninguna genera información adicional sobre el tráfico que no está en los datos.}
\end{frame}

\begin{frame}{Wavelet Haar: aproximación y detalle}
\small
Para bloques diádicos de tamaño \(s\), definimos la reconstrucción de la aproximación:
\[
A_s(t)=\frac{1}{s}\sum_{j\in\operatorname{bloque}_s(t)}x_j,
\qquad D_s(t)=A_{s/2}(t)-A_s(t).
\]
\begin{itemize}
\item \(A_s\) resume el promedio de demanda a escala \(s\).
\item \(D_s\) conserva el cambio que desaparece al aumentar la escala.
\item Se implementó la reconstrucción Haar en \textbf{NumPy} (sin dependencia PyWavelets).
\item Ejemplos: 2, 8 y 32 intervalos de 30 minutos equivalen a 1, 4 y 16 horas.
\end{itemize}
\end{frame}

\begin{frame}{Haar multiescala sobre la demanda}
\centering
\grafico{haar_multiescala.png}{.68\textheight}
\vspace{1mm}
{\scriptsize Señal de entrenamiento: la aproximación conserva la estructura general; el detalle destaca fluctuaciones a distintas escalas.}
\end{frame}

\begin{frame}{27 características de demanda Fourier/Haar}
\small
Se construyeron señales con inicios de viajes de \textbf{solo entrenamiento}, por región aproximada:
\begin{itemize}
\item \textbf{3 regiones}: general, Manhattan y alrededores de aeropuertos.
\item Por región: 6 aproximaciones/detalles Haar seleccionados y 3 componentes FFT.
\item Total: \(3\times(6+3)=\boxed{27}\) variables.
\end{itemize}
\begin{block}{Transformación estable}
\[
 x_t=\log\!\big(1+\mathrm{conteos}_t\cdot100000/n_{\mathrm{ref}}\big).
\]
\small Los conteos se ajustan al tamaño del conjunto de referencia y se suavizan con \(\log(1+\cdot)\).
\end{block}
\porq{La demanda se interpreta como una \textbf{señal proxy}, no como medición directa de velocidad ni de tráfico.}
\end{frame}

\begin{frame}{14 características espectrales basadas en duración}
\small
Para la media hora de la semana \(b\), se estima la duración logarítmica típica, suavizada:
\[
\widetilde\mu_b=\frac{\sum_{i\in b} \log(1+y_i)+45\overline z}{n_b+45}.
\]
\begin{itemize}
\item \textbf{5} reconstrucciones FFT conservando armónicas bajas.
\item \textbf{8} aproximaciones/detalles Haar a 4 escalas.
\item \textbf{1} media semanal suavizada.
\end{itemize}
\begin{block}{Condición indispensable}
\textcolor{Warning}{Estas 14 variables \textbf{sí utilizan el objetivo}}: deben generarse \textbf{out-of-fold} en train; para validación/test se ajustan solamente con el train permitido.
\end{block}
\end{frame}

\begin{frame}{Prueba OOF: la fila no puede verse a sí misma}
\small
\begin{enumerate}
\item Separar train en \textbf{5 folds} (semilla 47).
\item Para el fold \(j\), construir promedios y wavelets a partir de los \textbf{otros cuatro folds}.
\item Transformar únicamente los registros del fold \(j\).
\item Concatenar los cinco resultados en el orden original.
\item Ajustar un codificador con \textbf{todo el train permitido} para transformar validación/test.
\end{enumerate}
\begin{block}{Prueba programática documentada}
Al cambiar \(y_i\) de una observación, sus propias columnas supervisadas OOF \textbf{no se modifican}. Es una verificación parcial, no una garantía formal de ausencia de toda fuga.
\end{block}
\end{frame}

\begin{frame}{Inventario preciso de características}
\small
\begin{center}
\begin{tabular}{lrc}
\toprule
\textbf{Origen} & \textbf{Cantidad nueva} & \textbf{Acumulado}\\
\midrule
Originales y geográficas básicas & 81 & 81\\
Fourier/Haar deterministas & 45 & 126\\
Demanda regional en varias escalas & 27 & 153\\
Estadísticas espectrales de duración OOF & 14 & \textbf{167}\\
\bottomrule
\end{tabular}
\end{center}
\begin{block}{Tres configuraciones efectivamente entrenadas}
\centering \textbf{126} (bases fijas) $\longrightarrow$ \textbf{153} (+demanda) $\longrightarrow$ \textbf{167} (+OOF).
\end{block}
\porq{Esta comparación por etapas permite saber \emph{qué familia de variables aporta información} y cuál podría generar ruido.}
\end{frame}

\section{Experimentos espectrales y selección}
\begin{frame}{Modelo espectral: ¿por qué LightGBM?}
\small
\begin{itemize}
\item Aprende interacciones no lineales entre \textbf{distancia, hora, origen, destino y periodicidad}.
\item Tolera bien relaciones con umbrales y escalas diferentes, sin \texttt{StandardScaler} obligatorio.
\item Permite regularización y control de complejidad con muchos registros.
\item El modelo se entrena con \(z=\log(1+\texttt{trip\_duration})\).
\end{itemize}
\begin{block}{Configuración comparable}
\small \texttt{n\_estimators=1700; learning\_rate=0.045; num\_leaves=105;}\\
\texttt{min\_child\_samples=75; reg\_lambda=5; colsample=0.92;}\\
\texttt{subsample=0.90; seed=42; early\_stopping=110} en validación.
\end{block}
\end{frame}

\begin{frame}{Ablación: no todas las wavelets mejoran}
\small
\begin{center}
\begin{tabular}{lrc}
\toprule
\textbf{Modelo individual} & \textbf{Variables} & \textbf{RMSLE val.}\\
\midrule
Referencia original & 81 & 0.375798\\
Fourier/Haar fijas & 126 & 0.376372\\
Fourier + demanda Haar & 153 & 0.376802\\
Fourier + demanda + target OOF & 167 & 0.376833\\
\bottomrule
\end{tabular}
\end{center}
\begin{block}{Interpretación}
Las variantes espectrales \textbf{solas} fueron peores que la referencia en esta partición. Sin embargo, pueden cometer errores distintos y resultar útiles en combinación.
\end{block}
\porq{No se seleccionó el modelo espectral individual como ganador por tener más variables.}
\end{frame}

\begin{frame}{¿Por qué probar un ensamble y en qué escala?}
\small
Dos modelos complementarios pueden compensar errores. La combinación se hace \textbf{en log-duración}, la escala alineada con RMSLE:
\[
\widehat z = w\,\widehat z_{\mathrm{geo}}+(1-w)\widehat z_{\mathrm{espectral}},
\qquad \widehat y=\max\!\left(0,e^{\widehat z}-1\right).
\]
\begin{itemize}
\item Se compararon los pesos \(w\) en \textbf{validación}.
\item Se conservaron las predicciones del mejor ensamble geográfico anterior.
\item Se evaluaron las 3 variantes espectrales y sus combinaciones.
\end{itemize}
\begin{block}{Regla}Nunca ajustar los pesos usando holdout ni \texttt{test.csv}.\end{block}
\end{frame}

\begin{frame}{Comparación final en validación}
\centering
\grafico{comparacion_modelos_wavelet.png}{.65\textheight}
\vspace{1mm}
{\scriptsize Menor es mejor. Se mantuvieron las variantes que empeoraron para documentar el proceso de selección.}
\end{frame}

\begin{frame}{Ganador: 75\% previo + 25\% espectral}
\begin{columns}[T]
\column{.55\textwidth}
\small
\begin{itemize}
\item \textbf{75\%}: ensamble geográfico anterior, entrenado sin fuga.
\item \textbf{25\%}: LightGBM de 167 variables con Fourier + Haar + target OOF.
\item Pesos elegidos \textbf{solo por validación}.
\item Predicciones combinadas en escala \(\log(1+y)\).
\end{itemize}
\column{.43\textwidth}
\begin{block}{Resultado de validación}\centering
{\Large \textbf{0.373800}}\\[1mm] RMSLE
\end{block}
\begin{block}{Previo}\centering
{\large 0.374624}\\[-1mm] RMSLE
\end{block}
\end{columns}
\porq{La reducción fue de aproximadamente \(0.000824\) RMSLE en validación; es pequeña pero observable.}
\end{frame}

\begin{frame}{Prueba sellada: resultados de holdout}
\small
\begin{center}
\begin{tabular}{lcc}
\toprule
\textbf{Evaluación} & \textbf{Geográfico previo} & \textbf{Fourier + Haar}\\
\midrule
Validación (70,000) & 0.374624 & \textbf{0.373800}\\
Holdout (70,000) & 0.379288 & \textbf{0.378590}\\
\bottomrule
\end{tabular}
\end{center}
\vspace{2mm}
\begin{block}{Disciplina de evaluación}
\begin{itemize}
\item La elección del peso se hizo antes de consultar el holdout.
\item El conjunto ciego de 100,000 viajes \textbf{nunca} aporta etiquetas.
\item El puntaje externo de \textbf{0.3704} fue reportado por el usuario para otro envío y \textbf{no es comparable directamente} con estos valores locales.
\end{itemize}
\end{block}
\end{frame}

\begin{frame}{Diagnóstico: los viajes extremos dominan parte del error}
\small
\begin{center}
\begin{tabular}{lrr}
\toprule
\textbf{Duración real (validación)} & \textbf{Casos} & \textbf{Porción del error log-cuadrático}\\
\midrule
1 a 60 s & 433 & 10.21\%\\
61 a 300 s & 10,149 & 14.04\%\\
301 a 900 s & 35,997 & 26.33\%\\
901 a 3,600 s & 22,827 & 22.63\%\\
3,601 a 43,200 s & 491 & 2.38\%\\
Mayor de 12 h & \textbf{103} & \textbf{24.41\%}\\
\bottomrule
\end{tabular}
\end{center}
\porq{No se eliminaron extremos de validación: hacerlo produciría una estimación artificialmente favorable. Estos 103 casos son un diagnóstico, no un límite teórico para RMSLE.}
\end{frame}

\begin{frame}{Limitación metodológica: evaluación temporal}
\small
\begin{block}{Qué sí garantiza el experimento local}
\begin{itemize}
\item No se usan etiquetas del \texttt{test.csv}.
\item La mediana y las señales supervisadas se ajustan solo con train.
\item OOF impide incorporar la propia etiqueta en los predictores de esa fila.
\end{itemize}
\end{block}
\begin{alertblock}{Qué \emph{no} garantiza}
La partición aleatoria puede incorporar información de etiquetas de \textbf{fechas futuras} dentro del entrenamiento de otra fecha. Para pronosticar viajes futuros en producción se requerirían validación cronológica y agregaciones estrictamente causales.
\end{alertblock}
\end{frame}

\section{Reentrenamiento, entrega y conclusión}
\begin{frame}{Entrenamiento definitivo después de seleccionar}
\small
\begin{enumerate}
\item Tomar \textbf{700,000} filas etiquetadas y filtrar \textbf{solo en el ajuste} duraciones fuera de \([1,43200]\) s.
\item Quedan \textbf{699,016} filas para reajustar el componente final.
\item Reconstruir OOF, ajustar medianas \textbf{solo sobre entrenamiento} y entrenar el LightGBM de \textbf{167 variables}.
\item Combinar \(75/25\) con el modelo geográfico anterior, ya reentrenado.
\item Predecir cada registro de \texttt{test.csv} y transformar con \(\exp(z)-1\), recortando a cero.
\end{enumerate}
\begin{block}{Sin contaminación del test}La predicción no contiene ninguna medición de \texttt{trip\_duration} del conjunto ciego: no existe en \texttt{test.csv}.\end{block}
\end{frame}

\begin{frame}[fragile]{Verificación del CSV entregable}
\small
Formato requerido, \textbf{sin columnas auxiliares}:
\begin{lstlisting}[language=Python,basicstyle=\ttfamily\scriptsize]
submission = pd.DataFrame({
    'id': test['id'], 'trip_duration': pred
})
assert submission.columns.tolist() == ['id','trip_duration']
assert len(submission) == 100_000
assert submission['id'].is_unique
assert submission['id'].equals(test['id'])
assert np.isfinite(submission['trip_duration']).all()
assert (submission['trip_duration'] >= 0).all()
submission.to_csv('submission_wavelet_fft.csv', index=False)
\end{lstlisting}
\porq{La calidad de un modelo también depende de que el archivo final cumpla exactamente las condiciones técnicas de la evaluación.}
\end{frame}

\begin{frame}{Mapa de los archivos y reproducción}
\small
\begin{tabular}{p{.36\textwidth}p{.55\textwidth}}
\toprule
\textbf{Archivo / módulo} & \textbf{Para qué sirve}\\
\midrule
Notebook principal (\texttt{.ipynb}) & Fases, fórmulas, gráficos y evaluación\\
\texttt{features\_taxi.py} & Construye 81 variables base\\
\texttt{wavelet\_fft\_taxi.py} & Fourier, Haar y codificador OOF\\
Script de experimentos & Entrena tres variantes y guarda métricas\\
Script de ensambles & Selección exclusiva en validación\\
Script de holdout & Evaluación final reservada\\
Script de reentrenamiento & Crea el CSV con el modelo final\\
Script de predicción & Reproduce la entrega desde modelos guardados\\
\bottomrule
\end{tabular}
\vspace{2mm}
{\scriptsize Los nombres completos y el orden de ejecución están en \texttt{README\_Wavelets\_Fourier.md}.}
\end{frame}

\begin{frame}{Conclusiones: qué aprendimos del experimento}
\small
\begin{enumerate}
\item El reto es \textbf{regresión}: RMSLE es la medida objetivo, no F1 ni accuracy.
\item La ingeniería de características suele importar más que probar algoritmos al azar.
\item \textbf{Fourier} detectó ciclos temporales; \textbf{Haar} representó variaciones multiescala.
\item Las variantes espectrales \textbf{individuales} empeoraron; el \textbf{ensamble} mejoró ligeramente.
\item El mejor holdout pasó de \textbf{0.379288} a \textbf{0.378590}, con selección por validación.
\item El objetivo \textbf{0.20 no se consiguió}; la puntuación ciega del nuevo CSV aún debe medirse externamente.
\end{enumerate}
\end{frame}

\begin{frame}{Cómo defender el proyecto en preguntas}
\small
\begin{description}
\item[¿Por qué no F1?] Porque la variable objetivo es continua, no una clase.
\item[¿Por qué logaritmo?] Para alinear el entrenamiento con RMSLE y controlar errores multiplicativos.
\item[¿Por qué LightGBM?] Captura interacciones no lineales, escala a muchos viajes y no exige estandarización.
\item[¿Por qué OOF?] Para que los promedios de \texttt{trip\_duration} no incluyan el objetivo de la fila que predicen.
\item[¿Por qué mezclar?] Los errores de un modelo geométrico y uno espectral pueden ser complementarios.
\item[¿Llegamos a 0.20?] No. Reportamos valores medidos y sus límites.
\end{description}
\end{frame}

\begin{frame}{Referencias y material de apoyo}
\footnotesize
\begin{itemize}
\item \textbf{U5 -- 01 -- ML.pdf}, \emph{Introducción al Machine Learning}: aprendizaje supervisado, regresión y función de costo.
\item \textbf{U5 -- 02 -- Perceptrón.pdf}: descenso de gradiente y distinción entre activaciones de regresión/clasificación.
\item \textbf{03 -- Ciclo de vida del ML.pdf}: CRISP-DM, preprocesamiento, partición y prevención de \emph{data leakage}.
\item Scripts y notebook del proyecto, incluidos en este ZIP: evidencia de implementación y resultados, no resultados hipotéticos.
\end{itemize}
\vspace{2mm}
\centering\textcolor{UNAMteal}{\Large \textbf{Gracias. ¿Preguntas?}}
\end{frame}

\end{document}
